\documentclass[11pt,a4paper]{article}
\usepackage[utf8]{inputenc}
\usepackage[margin=1in]{geometry}
\usepackage{graphicx}
\usepackage{hyperref}
\usepackage{booktabs}
\usepackage{xcolor}
\usepackage{cite}
\usepackage{subcaption}

% tighter spacing to fit 5-6 pages
\setlength{\textfloatsep}{6pt plus 1pt minus 1pt}
\setlength{\floatsep}{6pt plus 1pt minus 1pt}
\setlength{\intextsep}{6pt plus 1pt minus 1pt}
\setlength{\abovecaptionskip}{2pt}
\setlength{\belowcaptionskip}{0pt}

% red marker for stuff I still need to fill in once training finishes
\newcommand{\todo}[1]{{\color{red}\textbf{[TODO: #1]}}}

% =====================================================================
% Planned figures (all saved to /kaggle/working/ by the notebook):
%   §4 Experimental setup
%     - split_overview.png  : bar chart of #clips per split, with #signers
%                             annotated; one bar group per split (train/val/test).
%                             Purpose: visualise that signer sets are disjoint.
%   §5 Results
%     - confusion_matrix.png: 100x100 confusion matrix of RQE-Transformer
%                             on signer-disjoint test, seed 0, normalised by row.
%     - training_curve.png  : train loss + val Top-1 vs epoch, RQE-Tx seed 0,
%                             single plot with twin y-axes.
%   §6 Failure analysis
%     - per_signer.png      : bar chart of test-set Top-1 per held-out signer,
%                             sorted ascending, with dashed line at macro mean.
%     - failure_example.png : 4 keyframes from a worst-signer misclassified
%                             clip with MediaPipe skeleton overlay; caption
%                             gives predicted vs true gloss.
% =====================================================================

\title{\vspace{-1.2cm}\textbf{Cross-Signer Generalization on WLASL100}\\
\large A pose-based RQE-Transformer evaluated on a signer-disjoint split}
\author{Patricia Guadalupe Alvarenga Mairena\\
Computer Vision II -- Master in Artificial Intelligence\\
Universidade de Vigo (ESEI)}
\date{\today}

\begin{document}
\maketitle

\begin{abstract}
Most word-level sign language recognition papers train and test on overlapping signer pools, so part of their accuracy is actually ``I remember this person''. I look at WLASL100 with a strict signer-disjoint split (no signer appears in both train and test) and ask whether a small Transformer over RQE-normalised MediaPipe landmarks helps on unseen signers, compared with a pooled-MLP that sees the same input but no temporal structure. Both models are trained from scratch under the same recipe and reported over 3 seeds.
\end{abstract}

%---------------------------------------------------------------
\section{Problem statement}

\textbf{Open challenge from the survey.} The reference survey for Domain~C (Zhou \emph{et al.}, 2023~\cite{zhou2023}) lists \emph{cross-subject / cross-signer generalization} as one of the unsolved problems in pose-based action recognition: models trained on one set of subjects tend to overfit to body proportions, signing style and camera placement, and drop noticeably on people they were not trained on. The SLR survey by De Coster \emph{et al.}~\cite{decoster2023survey} flags the same gap. On WLASL~\cite{li2020}, which has few examples per gloss and a non-uniform distribution of signers per class, this is the dominant generalization gap.

\textbf{Why it matters.} A sign language recognizer that only works for the signers it was trained on is not useful in practice (accessibility, HCI). The signer-disjoint gap is basically what tells you whether the model learned the sign or learned the person.

\textbf{Research question.} On a strict signer-disjoint split of WLASL100, does a small Transformer over RQE-normalised MediaPipe Holistic landmarks beat a pooled-MLP baseline that consumes the same RQE features?

\textbf{Initial hypothesis.} Partially yes. RQE subtracts the shoulder midpoint and divides by the inter-shoulder distance, so absolute position and absolute body size are gone before the network sees the input. Self-attention should then weight the discriminative hand-shape transitions over the relatively stable body pose. I expect the Transformer to beat the pooled-MLP in macro F1 by at least $\sim$5\,pp. I do \emph{not} expect the gap to close: RQE does not normalise handedness, signing speed or finger anatomy.

%---------------------------------------------------------------
\section{Related work}

WLASL was introduced by Li \emph{et al.}~\cite{li2020} together with the 100/300/1000/2000 subsets and per-instance \texttt{signer\_id} field in \texttt{WLASL\_v0.3.json}, which I use to build the signer-disjoint split. Their strongest baselines were RGB-based (I3D~\cite{carreira2017i3d}) and quietly inherit signer identity through face / clothing / background; the gap to unseen signers is documented but not closed.

Pose-based methods take the opposite route. ST-GCN~\cite{yan2018} hard-codes a skeletal adjacency, awkward when fingers carry most of the discriminative signal. SPTN / SPOTER~\cite{bohacek2022sptn} drops the graph and runs a per-frame Transformer over MediaPipe landmarks, with reported gains over ST-GCN on WLASL100; this is the closest published architecture to mine. OpenHands~\cite{selvaraj2022openhands} releases pose-based baselines and flags signer-disjoint evaluation as an open problem. The gap that motivates this work: published pose-based numbers are almost always on signer-overlapping splits, so the headline accuracy hides how much is identity rather than sign. I close that loophole by re-splitting WLASL100 strictly by signer. Standard ingredients: Transformer encoder~\cite{vaswani2017attention}, AdamW~\cite{loshchilov2019adamw}, MediaPipe Holistic~\cite{lugaresi2019mediapipe}.

%---------------------------------------------------------------
\section{Architecture selection}
\label{sec:arch}

\textbf{Chosen: RQE-Transformer.}
Per frame I take 75 MediaPipe Holistic joints (33 body + 21+21 hand) in 3D, apply RQE (centre on shoulder midpoint, divide by inter-shoulder distance, clip to $\pm 20$, fallback scale 1.0 if torso is degenerate), and flatten to a 225-d vector. Then linear projection to $d_\text{model}=256$, learned positional embedding (max length 64), 4-layer / 8-head pre-norm Transformer (FFN 512, dropout 0.1, GELU), masked temporal mean pool, linear head $\to$ 100 logits.

\textbf{Why this should help cross-signer generalization.} RQE at the input kills the two cleanest identity signals (position, body size) before the network sees them — a representational prior, not something the model has to learn from a small training set. Self-attention then weights hand-shape transitions regardless of when they occur, which matters because signers go at different speeds. No fixed skeletal graph (unlike ST-GCN) means no signer-specific corrections baked into GCN weights.

\textbf{Baseline (Pooled-MLP).} Same RQE input, masked mean over time, then a 2-layer MLP. Isolates how much the temporal model adds on top of RQE.

\textbf{Rejected alternatives.} ST-GCN~\cite{yan2018}: fixed graph hurts on finger-dominated signal. Raw RGB (I3D~\cite{carreira2017i3d}, Video-Swin): carries the strongest identity cues; not the comparison I want.

%---------------------------------------------------------------
\section{Experimental setup}
\label{sec:setup}

\textbf{Dataset.} WLASL100 (top-100 glosses). The chinhde Kaggle landmarks dump turned out to be incompatible with the official WLASL metadata (no \texttt{video\_id}, mismatched per-class counts), so I extracted landmarks directly from the raw videos in \texttt{risangbaskoro/wlasl-processed} using MediaPipe Holistic~\cite{lugaresi2019mediapipe}. Of the 2{,}038 WLASL100 video IDs, \todo{N\_present}/2038 were actually on disk; the rest have dead source URLs (see Limitations).

\textbf{Signer-disjoint protocol.} The official WLASL split is \emph{not} signer-disjoint: the same person can appear in train and test. I rebuild a split from \texttt{WLASL\_v0.3.json}'s \texttt{signer\_id} field, assigning whole signers (not videos) to splits: about 70\% train / 10\% val / 20\% test, with a greedy balancing pass so every class still has $\geq 1$ training clip. Final sizes: \todo{$N_\text{train}/N_\text{val}/N_\text{test}$ clips, $S_\text{train}/S_\text{val}/S_\text{test}$ signers}. The intersection of signer sets across splits is asserted empty.

\textbf{Metrics.} Top-1 accuracy and macro F1 on the signer-disjoint test set (held-out signers, never seen at training time). Per-signer accuracy is also reported as a diagnostic for the failure analysis.

\textbf{Training configuration.} AdamW lr $3\!\times\!10^{-4}$, weight decay $1\!\times\!10^{-4}$, cosine schedule, batch size 32, max 64 frames per clip, 15 epochs, cross-entropy with label smoothing 0.1, grad-clip 1.0. Trained from scratch (no pretraining). Each model is trained with 3 random seeds (0, 1, 2); I report mean $\pm$ std on the held-out test set.

\begin{figure}[t]
\centering
\includegraphics[width=0.5\linewidth]{split_overview.png}
\caption{Signer-disjoint split: \#clips per split with \#signers in parentheses. Signer sets are disjoint by construction.}
\label{fig:split}
\end{figure}

%---------------------------------------------------------------
\section{Results}

\begin{table}[h]
\centering
\caption{WLASL100, signer-disjoint test set. Mean $\pm$ std over 3 seeds.}
\label{tab:results}
\begin{tabular}{lcc}
\toprule
Model & Top-1 (\%) & Macro F1 (\%) \\
\midrule
Pooled-MLP (RQE input, no temporal)         & \todo{xx.x $\pm$ x.x} & \todo{xx.x $\pm$ x.x} \\
RQE-Transformer (this work)                 & \todo{xx.x $\pm$ x.x} & \todo{xx.x $\pm$ x.x} \\
\bottomrule
\end{tabular}
\end{table}

\begin{figure}[t]
\centering
\begin{subfigure}{0.48\linewidth}
  \includegraphics[width=\linewidth]{confusion_matrix.png}
  \caption{Confusion matrix (seed 0).}
  \label{fig:cm}
\end{subfigure}\hfill
\begin{subfigure}{0.48\linewidth}
  \includegraphics[width=\linewidth]{training_curve.png}
  \caption{Train loss / val Top-1 over 15 epochs (seed 0).}
  \label{fig:loss}
\end{subfigure}
\caption{RQE-Transformer on the signer-disjoint test set.}
\end{figure}

\textbf{Comparison with published numbers.} Pose-based methods with pretraining reach 60--80\% Top-1 on the \emph{standard} WLASL100 split~\cite{selvaraj2022openhands,bohacek2022sptn}. My setup is intentionally weaker (no pretraining, no augmentation, no hand-crops, half the videos missing, signer-disjoint split is harder), so I expect to be well below that. The point is the internal comparison between the two rows in Table~\ref{tab:results}, not the absolute number.

%---------------------------------------------------------------
\section{Failure analysis}
\label{sec:fail}

\textbf{Case 1 — Worst-performing held-out signer.}
Worst signer: \todo{signer\_id} with \todo{xx\%} Top-1 vs.\ macro mean \todo{yy\%} (Figure~\ref{fig:persigner}). Looking at the misclassified clips (Figure~\ref{fig:fail}), this signer is \todo{left-handed / unusually fast / camera-tilted}. RQE normalises position and torso scale but not handedness or speed, so the same gloss lands in a different region of \texttt{input\_proj}'s output, and the learned \emph{absolute} positional embedding misaligns attention patterns calibrated on slower clips. Both failures match design choices: the RQE invariance set is incomplete and the positional embedding is absolute rather than relative.

\begin{figure}[t]
\centering
\includegraphics[width=0.6\linewidth]{per_signer.png}
\caption{Per-signer Top-1 on the signer-disjoint test set (RQE-Tx). Dashed line: macro mean.}
\label{fig:persigner}
\end{figure}

\textbf{Case 2 — Top confused gloss pair.}
The biggest off-diagonal cell is \todo{gloss A vs gloss B} (\todo{xx\%} of A misclassified as B). The two glosses share the same gross arm trajectory and differ in hand-shape. MediaPipe hand landmarks are in the input, but fingers contribute $\sim$$\tfrac{42 \times 3}{225}\!\approx\!56\%$ of the per-frame vector \emph{by count} and yet much less than that in \emph{magnitude} after the linear projection (most variance is in the body channels). After the temporal mean pool, the discriminative finger signal gets averaged out. This is a representational limitation of pose-only input at MediaPipe's resolution, not a learning failure — ST-GCN would have the same problem.

\begin{figure}[t]
\centering
\includegraphics[width=0.85\linewidth]{failure_example.png}
\caption{Four keyframes of a misclassified test clip with skeleton overlay. Predicted \todo{X}, true \todo{Y}.}
\label{fig:fail}
\end{figure}

%---------------------------------------------------------------
\section{Thesis direction}

\textbf{Did the results support the hypothesis?}
\todo{Fill in once training finishes. Template: ``The RQE-Transformer beat the pooled-MLP by $\Delta$\,pp in macro F1 on the signer-disjoint test set, supporting the direction of the hypothesis. The gap is smaller than the $\sim$5\,pp I expected and is not uniform across signers (Section~\ref{sec:fail}), so the hypothesis is supported in direction but only partially in magnitude.''}

\textbf{What the failure analysis reveals.} Both failure cases trace to specific design choices, not generic difficulty. Case 1 (worst held-out signer) reflects an \emph{incomplete invariance set} in the input encoding: RQE removes position and torso scale but does nothing about handedness or signing speed, and the learned absolute positional embedding amplifies the speed mismatch. Case 2 (top confused gloss pair) reflects \emph{magnitude imbalance in the per-frame token}: the fine-grained finger signal is averaged out against the much stronger body-pose signal during the temporal mean pool. What needs to change is therefore the \emph{input representation and the pooling}, not necessarily a bigger Transformer.

\textbf{Concrete next step (proposed method).} Add a \emph{finger-resolution branch} to the architecture. Concretely: keep the body stream unchanged (75-joint RQE tokens through the current Transformer), and introduce a second token stream with only the 21 landmarks of the dominant hand (auto-detected from per-frame hand-motion energy and mirrored to a canonical side, which also addresses Case 1's handedness leak). The hand stream has its own learned positional embedding and its own short Transformer, and is fused into the body stream via cross-attention before the temporal mean pool.

\textbf{Feasibility.} All ingredients are preprocessing-only or additive on top of the current code: the hand landmarks are already in the cached \texttt{wlasl100\_landmarks.npz}; dominant-hand detection is a single pass over hand-motion energy; the extra Transformer adds $<$200k parameters relative to the current $\sim$3.2\,M, well within the same T4 training budget; evaluation is on the same signer-disjoint split, so the comparison is direct and the effect on Case 2 (and Case 1, via the canonicalisation) is directly measurable.

\textbf{Limitations.} About half of the WLASL100 video URLs are no longer reachable, so the test set is a subset of the benchmark and absolute numbers cannot be compared to public leaderboards (the internal comparison between the two rows of Table~\ref{tab:results} is still fair because both share the same input and recipe). Input is pose-only, so MediaPipe failures (zero hand landmarks on some frames) propagate into the model. Variance is estimated from only 3 seeds.

%---------------------------------------------------------------
\section*{Declaration of LLM use (per \S10 Academic Integrity)}

This project used an LLM coding assistant (GitHub Copilot Chat, Claude models). \emph{Drafted by the LLM, reviewed by me:} most boilerplate Python (PyTorch dataset, training loop, plotting helpers, MediaPipe extraction), notebook markdown, and the structural skeleton of this report. \emph{Drafted by me, polished by the LLM:} the research question, the choice of a strict signer-disjoint split, the mechanistic justifications (\S\ref{sec:arch}), the failure-analysis interpretations (\S\ref{sec:fail}), and the next-step proposal. \emph{Verified by me:} every number traced to a Kaggle cell output; citations checked manually; the signer-disjoint assertion is in the notebook. \emph{Not used:} I did not let the assistant invent citations or numerical results. One concrete catch: an earlier notebook used \texttt{chinhde/wlasl-300-landmarks} assuming it followed the official WLASL splits — it doesn't (no \texttt{video\_id}, mismatched per-class counts). A fingerprint-matching diagnostic confirmed the incompatibility and I rebuilt the pipeline on raw videos.

%---------------------------------------------------------------
\bibliographystyle{IEEEtran}
\begin{thebibliography}{10}
\bibitem{li2020}
D.~Li, C.~Rodriguez-Opazo, X.~Yu, and H.~Li, ``Word-level deep sign language recognition from video: A new large-scale dataset and methods comparison,'' in \emph{WACV}, 2020.

\bibitem{vaswani2017attention}
A.~Vaswani \emph{et al.}, ``Attention is all you need,'' in \emph{NeurIPS}, 2017.

\bibitem{lugaresi2019mediapipe}
C.~Lugaresi \emph{et al.}, ``MediaPipe: A framework for building perception pipelines,'' \emph{arXiv:1906.08172}, 2019.

\bibitem{yan2018}
S.~Yan, Y.~Xiong, and D.~Lin, ``Spatial temporal graph convolutional networks for skeleton-based action recognition,'' in \emph{AAAI}, 2018.

\bibitem{loshchilov2019adamw}
I.~Loshchilov and F.~Hutter, ``Decoupled weight decay regularization,'' in \emph{ICLR}, 2019.

\bibitem{selvaraj2022openhands}
P.~Selvaraj, G.~NC, P.~Kumar, and M.~M.~Khapra, ``OpenHands: Making sign language recognition accessible with pose-based pretrained models across languages,'' in \emph{ACL}, 2022.

\bibitem{bohacek2022sptn}
M.~Bohacek and M.~Hr\'uz, ``Sign pose-based transformer for word-level sign language recognition,'' in \emph{WACV Workshops}, 2022.

\bibitem{decoster2023survey}
M.~De Coster, D.~Shterionov, M.~Van Herreweghe, and J.~Dambre, ``Machine translation from signed to spoken languages: state of the art and challenges,'' \emph{Universal Access in the Information Society}, 2023.

\bibitem{zhou2023}
L.~Zhou \emph{et al.}, ``Human pose-based estimation, tracking and action recognition with deep learning: a survey,'' \emph{arXiv:2310.13039}, 2023.

\bibitem{carreira2017i3d}
J.~Carreira and A.~Zisserman, ``Quo vadis, action recognition? A new model and the Kinetics dataset,'' in \emph{CVPR}, 2017.
\end{thebibliography}

\end{document}
